\begin{abstract}
\noindent Retirement policies change the return to work decades before benefits are paid, act differently on workers in good and poor health, and spill over to income taxes, disability insurance, and means-tested transfers. We compute the marginal value of public funds (MVPF) of 66 retirement reforms in a life-cycle model of labor supply, saving, and Social Security claiming with health, disability, mortality, and wage risk, estimated on US panel data. For the 1920--1940 birth cohorts, the ranking turns on whether a reform rewards later claiming or lowers benefits at every claiming age. Raising the delayed retirement credit from the cohort-average 4.5 percent, or steepening the adjustment of benefits around the normal retirement age, raises welfare and lowers net government spending: these reforms pay for themselves. Raising the normal retirement age saves money but lowers welfare. Among benefit expansions, raising the bottom replacement rate of the benefit formula delivers the most welfare per dollar, and expanding Medicare coverage the least.

\medskip
\noindent\textbf{JEL codes:} H55, J26, H21, I13, D15.

\noindent\textbf{Keywords:} Social Security, retirement, marginal value of public funds, health, delayed retirement credit, life-cycle model.
\end{abstract}
